% Sections 13-16: Comparison with Existing Methods, Conclusion, Future Scope. Shared by both templates.

\section{Comparison with Existing Methods}\label{sec:compare}
Table~\ref{tab:compare} positions this work against the closest existing systems. Direct speed
comparisons are not meaningful, because the systems simulate different models on different
hardware; we compare what each one parallelises, how it handles reproducibility, and what it
measures.

\begin{widetable}
\centering
\caption{This work compared with existing methods.}
\label{tab:compare}
\footnotesize
\begin{tabularx}{\textwidth}{@{}p{2.4cm} X X X X@{}}
\toprule
System & What runs in parallel & Matching & Reproducibility & Scaling analysis\\
\midrule
ABIDES / ABIDES-Gym~\cite{byrd2020abides,amrouni2021} & nothing (serial event loop) & serial, one exchange & seed only & none\\
JAX-LOB~\cite{frey2023} & independent books on GPU & serial within book, on GPU & not discussed & per-message timing\\
PyMarketSim~\cite{mascioli2024} & nothing & serial, one asset & not discussed & single-thread throughput\\
FLAME GPU 2~\cite{richmond2023} & agents on GPU (spatial models) & not applicable (bidding sub-models) & atomics non-deterministic & agent-count scaling\\
Repast4Py~\cite{collier2022} & spatial partitions over MPI & not applicable & seed; cross-rank equality not addressed & strong scaling\\
SIMCoV-GPU~\cite{leyba2024} & spatial domains over GPUs & not applicable & statistical only ($\approx$1\%) & strong and weak scaling\\
\midrule
\textbf{This work} & agents on GPU or CPU cores; instruments over threads and MPI ranks & serial within instrument, parallel across instruments & \textbf{bit-identical trade log} on all back ends & strong, weak, Amdahl fit, Karp--Flatt, roofline, crossover\\
\bottomrule
\end{tabularx}
\end{widetable}

\textbf{Same conclusion, measured.} JAX-LOB states that matching inside one book is
serial~\cite{frey2023}. Our results quantify the consequence end to end: matching is 92--96\% of a
CUDA tick, so an 85$\times$ faster kernel yields only 1.1--1.6$\times$ for the whole simulation.

\textbf{Different decomposition.} Spatial ABM frameworks~\cite{richmond2023,collier2022,leyba2024}
partition by space and exchange halos. We partition by instrument, which removes all cross-rank
conflicts; the only per-tick communication is one integer reduction. The cost moves to load
imbalance at the barrier (20--39\% of tick time with 2--4 ranks), the analogue of Repast4Py's
communication-bound regime.

\textbf{Stronger reproducibility.} SIMCoV-GPU matches CPU and GPU results only
statistically~\cite{leyba2024}, and OpenRAND makes only the random numbers
reproducible~\cite{khan2024}. Counter-based random numbers, integer-only decisions, a deterministic
arrival order and order-independent compaction together give a bit-identical trade log across 12
configurations, so a parallel bug shows up as a changed fingerprint.

\textbf{Measured threshold.} Like SIMCoV-GPU, we find the GPU pays off only above a problem size;
we locate it at 4{,}000--16{,}000 agents against the 2-thread CPU and explain both the lower bound
(fixed copy and launch cost) and why the advantage shrinks again (growing matching time).

\textbf{Limits of this comparison.} Our CPU baseline on the T4 has only 2 threads; a stronger CPU
would move the crossover up. Our agents are simpler than ABIDES agents (no inventory or latency
model), and the model fails the no-autocorrelation stylised fact.

\section{Conclusion}\label{sec:concl}
We built a heterogeneous agent-based limit order book simulator with agent decisions in CUDA,
matching parallelised across instruments with OpenMP, and instruments distributed with MPI. All
eight objectives of the synopsis were met. Every back end produces a bit-identical trade log for a
given seed, which made correctness checkable throughout. The measurements give one consistent
answer: order matching under price--time priority is the scaling bottleneck. The GPU makes agent
decisions about 85$\times$ faster, yet the whole tick improves by at most 1.41$\times$ over the
CPU baseline, because matching takes 92--96\% of the time; the measured 1.11$\times$ at 1M agents
sits close to the Amdahl bound of 1.16$\times$. Matching parallelises only across instruments
(0.90$\times$ with one, 3.97$\times$ with 128). The GPU overtakes the multicore CPU between 4{,}000
and 16{,}000 agents, below which copy and launch costs dominate. On the laptop, CPU scaling levels
off near 2$\times$ (effective serial fraction 0.49) because of shared power and memory limits, and
MPI overhead is dominated by load imbalance at the per-tick barrier. The simulated returns show
heavy tails and volatility clustering, with excess short-term mean reversion as a known limitation.

\section{Future Scope}\label{sec:future}
\begin{itemize}
\item \textbf{Faster matching.} Match independent price levels speculatively, or batch orders in
  periodic call auctions, which allow parallel matching within an instrument at the cost of
  changing market rules.
\item \textbf{Moving the host scatter to the GPU.} Keep agent state and intents on the device and
  send per-instrument order lists directly, removing the 3.2\,ms host scatter.
\item \textbf{Load balancing for MPI.} Assign instruments to ranks by measured activity instead of
  contiguous blocks, to cut barrier waiting; and run on a real multi-node cluster to measure network
  cost.
\item \textbf{More realistic agents.} Add inventory, cash and latency, and tune the agent mix to
  remove the negative lag-1 autocorrelation.
\item \textbf{Stronger hardware study.} Repeat the crossover on a server CPU with many cores and on
  newer GPUs, and confirm the hardware-saturation explanation with performance counters.
\item \textbf{Multi-GPU.} Give each MPI rank its own GPU for very large agent populations.
\end{itemize}
